\chapter{Simulação, teste e comissionamento}
\label{cap:simulacao}

\begin{objetivos}
Ao final deste capítulo você deve ser capaz de:
\begin{itemize}
  \item explicar por que testar antes de energizar é obrigatório;
  \item percorrer os seis níveis de validação, do teste de lógica ao \emph{dry-run};
  \item montar um roteiro de teste com casos, estímulos, resultados esperados e evidências;
  \item reconhecer os limites da simulação --- o que ela não detecta.
\end{itemize}
\end{objetivos}

\section{Por que testar antes}

Um controlador mal programado pode causar perda de produção e criar condições perigosas para
equipamentos e pessoas. Simular permite encontrar e corrigir defeitos cedo no projeto, quando a
correção custa minutos em vez de horas de parada. Além disso, melhora a segurança associada ao
equipamento e reduz o tempo de instalação e comissionamento, porque muitos cenários já foram
testados antes de o sistema ser energizado.

\index{simulação de programa}\index{teste de programa}

\section{A escada de validação}

A validação é uma escada: cada degrau é mais caro e mais realista que o anterior, e cada degrau
elimina uma classe de defeito. Subir todos os degraus antes de energizar a máquina é o que evita
surpresa no comissionamento.

\begin{figure}[H]
  \centering
  \includegraphics[width=\linewidth]{16-escada-validacao.pdf}
  \caption{Níveis de validação, do teste de lógica ao comissionamento.}
  \label{fig:escada-validacao}
\end{figure}

\begin{table}[H]
  \centering
  \caption{O que cada nível detecta e o que ele não detecta.}
  \label{tab:niveis-validacao}
  \begin{tabular}{p{3.6cm} p{5.4cm} p{5.4cm}}
    \toprule
    Nível & Detecta & Não detecta \\
    \midrule
    1. Revisão de lógica
      & Erro de lógica, endereço trocado, condição faltante, laço mal construído
      & Nada de comportamento em tempo real \\
    2. Simulação de E/S
      & Sequência, temporização relativa, intertravamento, tabela-verdade
      & Comportamento de sinal elétrico e de campo \\
    3. Controlador virtual
      & Portabilidade do código, desempenho relativo, integração com o ambiente
      & Diferenças de implementação do hardware real \\
    4. Bancada com E/S
      & Fiação de campo, tensões, sensores e atuadores reais, ruído
      & Comportamento da máquina sob carga \\
    5. HIL
      & Interação controlador--planta simulada, casos difíceis e de falha
      & Desgaste e efeitos mecânicos reais \\
    6. \emph{Dry-run} em campo
      & Tudo o que envolve a instalação completa, com o processo energeticamente seguro
      & Desempenho em produção plena \\
    \bottomrule
  \end{tabular}
\end{table}

\paragraph{Simulação de E/S.} É o nível que todo aluno tem disponível: o ambiente de programação
permite forçar o estado das entradas e observar as saídas sem nenhum hardware conectado. Com ela se
percorre a árvore de decisões do programa --- inclusive as combinações que a máquina real demora
horas para alcançar.

\paragraph{Controlador virtual (\emph{soft PLC}).} A função de controlador é executada em hardware de
propósito geral, em máquina virtual ou contêiner. Permite ensaiar o programa em outra plataforma e é
útil em ensino e em testes de integração.

\paragraph{HIL (\emph{hardware-in-the-loop}).} O controlador real roda conectado a um modelo
computacional da planta. É o nível que permite ensaiar falha de sensor, sobrecurso de válvula e
situação de emergência sem risco e sem parar a produção.

\paragraph{\emph{Dry-run} em campo.} Máquina energizada, sem produção: verifica-se fim de curso,
polaridade, tempo de resposta e cada intertravamento, ponto por ponto, com a equipe de campo
presente.

\section{Roteiro de teste}

Teste sem roteiro é passeio. O roteiro abaixo é um modelo para adaptação; a coluna de evidência é a
que costuma ser esquecida e é a que sustenta a entrega técnica.

\begin{table}[H]
  \centering
  \caption{Modelo de roteiro de teste.}
  \label{tab:roteiro-teste}
  \begin{tabular}{p{1.0cm} p{3.4cm} p{3.6cm} p{3.4cm} p{2.6cm}}
    \toprule
    \# & Caso & Estímulo & Resultado esperado & Evidência \\
    \midrule
    T01 & Partida em condições normais
      & Todas as condições de partida satisfeitas; comando de partida
      & Máquina parte na sequência prevista
      & Registro no relatório de comissionamento \\
    T02 & Bloqueio de partida por condição
      & Uma condição de partida não satisfeita; comando de partida
      & Máquina não parte e alarme indicado
      & Lista de alarmes ativos \\
    T03 & Intertravamento em operação
      & Simular perda da condição de intertravamento durante a operação
      & Parada segura na sequência prevista
      & Registro de eventos com horário \\
    T04 & Parada de emergência
      & Atuar o dispositivo de emergência
      & Parada conforme projeto; retomada só após desbloqueio e reset
      & Teste presencial com assinatura \\
    T05 & Retomada após falta de energia
      & Desenergizar e reenergizar o painel
      & Estado seguro; variáveis retentivas preservadas
      & Leitura das variáveis retentivas \\
    T06 & Falha de sensor
      & Desconectar sensor em bancada
      & Programa assume o comportamento definido no descritivo (e seguro)
      & Registro de falha e ação tomada \\
    T07 & Falha de comunicação
      & Interromper o enlace com a rede ou com a remota de E/S
      & Comportamento seguro declarado; alarme de comunicação
      & Registro de eventos \\
    T08 & Tempo de ciclo
      & Operação em carga máxima
      & Tempo de varredura medido dentro do requisito
      & Medição registrada no relatório \\
    \bottomrule
  \end{tabular}
\end{table}

Os casos T04 a T07 são os que a simulação isolada \textbf{não} cobre adequadamente, porque dependem
de comportamento de hardware, de energia ou de rede. Eles precisam ser testados com equipamento real
--- e sempre pelo procedimento de bloqueio da instalação.

\section{Limites da simulação}

\begin{table}[H]
  \centering
  \caption{O que a simulação não substitui.}
  \label{tab:limites-simulacao}
  \begin{tabular}{p{5.6cm} p{9.0cm}}
    \toprule
    Limite & Por quê \\
    \midrule
    Sinal elétrico real & Tensão, corrente, ruído, tempo de comutação e falha de isolamento só
      aparecem com módulo e carga reais \\
    Dinâmica mecânica & Inércia, folga, atrito e sobrecurso não existem no modelo \\
    Atraso de rede & Latência e perda de pacote dependem da instalação \\
    Falha do próprio controlador & Nenhuma simulação executa o programa com o controlador travado \\
    Erro de campo & Sensor instalado no lugar errado, inversão de polaridade e fiação trocada só
      aparecem na bancada e no \emph{dry-run} \\
    \bottomrule
  \end{tabular}
\end{table}

\begin{atencao}
Nunca energize máquina com o programa recém-carregado sem ter percorrido o roteiro de teste. E nunca
faça \emph{dry-run} com pessoa em zona de risco: a máquina está energizada, com lógica nova e sem
histórico --- a condição de maior risco de todo o projeto.
\end{atencao}

\section{Registro de evidências}

O resultado do teste precisa sair do ambiente e virar documento:

\begin{itemize}
  \item roteiro preenchido, com data, responsável e assinatura nos casos de segurança;
  \item registro de eventos com data e hora dos alarmes e paradas testadas;
  \item lista de pendências encontradas e o que foi feito com cada uma;
  \item versão do programa testada (número de revisão) --- teste sem versão identificada não vale
        como evidência.
\end{itemize}

Esse conjunto é o que liga este capítulo ao anterior: a documentação da
Capítulo~\ref{cap:documentacao} é o que permite provar, meses depois, que o sistema foi testado.

\begin{pratica}
Monte o roteiro de teste para um dos exercícios de laboratório desta apostila, preenchendo as colunas
da Tabela~\ref{tab:roteiro-teste}. Depois:
(a) execute os casos possíveis no simulador de E/S e registre o resultado obtido;
(b) identifique quais casos só podem ser executados com hardware real;
(c) escreva, em cinco linhas, o risco de liberar o programa sem executar esses casos.
\end{pratica}

\begin{exercicios}
\begin{enumerate}[label=\alph*)]
  \item Cite três benefícios de testar em simulação antes da instalação.
  \item Ordene os seis níveis de validação do menos para o mais realista e explique o que cada um
        acrescenta.
  \item Por que a simulação de E/S não substitui o teste com módulo de entrada real?
  \item O que é HIL e que tipo de falha ele permite ensaiar sem risco?
  \item Escreva dois casos de teste adicionais para uma máquina de envase, no formato da
        Tabela~\ref{tab:roteiro-teste}.
  \item Por que a versão do programa é parte obrigatória da evidência de teste?
  \item Um fornecedor entrega o sistema afirmando que ``já testou tudo no simulador''. Que
        questionamento você faz, com base neste capítulo?
  \item Explique por que o \emph{dry-run} é a etapa de maior risco do projeto.
\end{enumerate}
\end{exercicios}
